% Standalone summary of the two ChatGPT knowledge-design elicitation sessions.
% This document deliberately does not reproduce prompt or response text.
\documentclass[11pt,a4paper]{article}

\usepackage[margin=28mm]{geometry}
\usepackage[T1]{fontenc}
\usepackage[utf8]{inputenc}
\usepackage{lmodern}
\usepackage{microtype}
\usepackage{booktabs}
\usepackage{tabularx}
\usepackage{hyperref}
\usepackage{fancyhdr}

\hypersetup{
  pdftitle={CliniCause Knowledge-Design Stage: ChatGPT Elicitation Summaries},
  pdfauthor={CliniCause authors},
  pdfsubject={Both sessions used ChatGPT 5.4 with extended reasoning},
  pdfkeywords={CliniCause, ChatGPT 5.4, extended reasoning, MIMIC-III, PhysioNet, supplementary material},
  colorlinks=true, linkcolor=black, urlcolor=blue
}

\pagestyle{fancy}
\fancyhf{}
\fancyfoot[C]{\thepage}
\renewcommand{\headrulewidth}{0pt}
\renewcommand{\footrulewidth}{0pt}
\setlength{\parindent}{0pt}
\setlength{\parskip}{0.55em}
\newcommand{\sourcepath}[1]{{\footnotesize\nolinkurl{#1}}}
\newcommand{\mimichash}{{\scriptsize\ttfamily 56271465\allowbreak a73f6241\allowbreak 51ba36b1\allowbreak b4c4a07a\allowbreak d0e4271a\allowbreak 9bb73042\allowbreak 56ede415\allowbreak 5afbf79a}}
\newcommand{\physionethash}{{\scriptsize\ttfamily 933c864a\allowbreak 5c067976\allowbreak 4d579d46\allowbreak 17f5192c\allowbreak 7f1db0a5\allowbreak 3cd74e5b\allowbreak 0906b729\allowbreak 7d98cb67}}

\begin{document}

\hypersetup{pageanchor=false}
\begin{titlepage}
  \centering
  \vspace*{0.20\textheight}
  {\LARGE\bfseries Supplementary Material\par}
  \vspace{1.2cm}
  {\huge\bfseries CliniCause Knowledge-Design Stage\par}
  \vspace{0.35cm}
  {\Large ChatGPT Elicitation Summaries\par}
  \vfill
  \begin{minipage}{0.84\textwidth}
    \centering\large
    \textbf{Model metadata: ChatGPT 5.4 with extended reasoning}\\[0.4em]
    The same model setting was used for the MIMIC-III and PhysioNet sessions.
  \end{minipage}
  \vfill
  {\large Standalone academic supplement\par}
\end{titlepage}
\hypersetup{pageanchor=true}
\pagenumbering{arabic}

\tableofcontents
\clearpage

\section{Purpose, scope, and handling of elicitation output}
\textbf{Both elicitation sessions used ChatGPT 5.4 with extended reasoning.}

The CliniCause knowledge-design stage used ChatGPT to propose clinically and
causally motivated design candidates for two ICU datasets. The model received
dataset schema/column names and the mortality outcome, but no patient rows,
empirical distributions, prevalence values, or causal effect estimates. It was
asked to propose proxy variables, clinical cutoffs, temporal and aggregation
rules, supporting literature, and dataset-specific directed acyclic graph (DAG)
relationships.

The elicitation was an input to author review rather than an automated labeling
or analysis step. The authors inspected the proposals and implemented the
proposals they accepted. This supplement summarizes the records; it does not
copy, quote, reproduce, correct, or complete any prompt or model response.

\begin{table}[ht]
\centering
\small
\begin{tabularx}{\textwidth}{@{}lXX@{}}
\toprule
Session & Source record & Record scope summarized here \\
\midrule
MIMIC-III & \sourcepath{thesis-writing/prompts-and-documents/mimic-prompt-running.pdf} & Latent-variable causal design using the MIMIC-III schema and mortality outcome. \\
PhysioNet & \sourcepath{thesis-writing/prompts-and-documents/physionet-prompt-running.pdf} & Latent-variable causal design using the PhysioNet ICU schema and in-hospital mortality outcome. \\
\bottomrule
\end{tabularx}
\end{table}

\section{MIMIC-III session summary}
\textbf{Source-file provenance.} \sourcepath{thesis-writing/prompts-and-documents/mimic-prompt-running.pdf}\\
\textbf{Source PDF title.} \emph{mimic-prompt-running - Latent Causal Model MIMIC-III}\\
\textbf{Source length and checksum.} 44 pages; SHA-256 \mimichash.

The MIMIC-III session established a staged workflow: audit the available
tables/variables; identify background, measurement, care-process, and outcome
variables; select latent clinical states; operationalize those states with
binary proxy rules; and then specify a sparse causal DAG. The response treated
observed laboratory values, vital signs, and charted measures primarily as
manifestations of latent physiology rather than as a dense network of direct
causes.

The proposed design covered acute physiologic and organ-dysfunction domains,
with attention to chronic baseline risk and global severity. It mapped observed
measurements to candidate latent states, supplied implementable scoring or
threshold-based proxy-label logic, and described missingness and measurement
intensity as potentially informative care-process variables. It also included
temporal safeguards against future information and outcome leakage, plus
guidance on distinguishing confounders, mediators, and colliders for analyses
of mortality.

The session further supplied candidate latent-to-latent, latent-to-observation,
and latent-to-mortality relationships; pseudocode and an edge-list-oriented
implementation representation; and literature-oriented rationale for clinical
constructs and causal assumptions. The record flags limitations where a latent
state is weakly measured or where care actions and documentation can complicate
interpretation.

\section{PhysioNet session summary}
\textbf{Source-file provenance.} \sourcepath{thesis-writing/prompts-and-documents/physionet-prompt-running.pdf}\\
\textbf{Source PDF title.} \emph{physionet-prompt-running - Causal Model for ICU Data}\\
\textbf{Source length and checksum.} 45 pages; SHA-256 \physionethash.

The PhysioNet session used the same latent-variable causal-design framework,
but adapted the operationalization to the variables available in that ICU
dataset. It audited demographic, vital-sign, laboratory, and care-related
fields; proposed candidate latent physiologic and organ-dysfunction states; and
defined binary proxy-label rules with clinically motivated thresholds and
aggregation over observed signals.

The response explicitly emphasized the dataset's information limits: several
clinical concepts cannot be directly identified when diagnoses, medications,
procedures, microbiology, comorbidity coding, admission diagnosis, or clinician
intent are absent. Accordingly, it presented states such as sepsis, shock, and
cardiac injury as proxy constructs rather than directly observed ground truth.
It also considered missingness and measurement intensity, temporal leakage,
and the role of care processes when interpreting measurements.

The proposed DAG separated background/case-mix factors, latent acute
physiology, observed measurements, care or measurement processes, and
in-hospital mortality. It included causal-analysis guidance on adjustment,
mediators, and collider risk; code-oriented representations for implementation;
and supporting clinical and causal literature directions. The session marked
the limits of available evidence and treated proposed label rules as
operational approximations for subsequent author review.

\section{Record-integrity notes}
The two source records are distinct PDFs with distinct titles, session links,
page counts, and checksums. Their requested staged methodology is closely
parallel, while the dataset audits and later operational details are
dataset-specific. No duplicate source file was used in preparing this summary.

Text extraction identifies a trailing blank page in each source record (page 44
for MIMIC-III and page 45 for PhysioNet); this appears to be a source-export
artifact, not missing prompt material. The source records themselves do not
provide patient-level data, prevalence summaries, distributions, or causal
effect estimates. Model-setting metadata is stated in this supplement as
provided for the study: ChatGPT 5.4 with extended reasoning.

\end{document}
